\documentclass[11pt,a4paper]{article}
\usepackage[T1]{fontenc}
\usepackage{isabelle,isabellesym}

\usepackage{amsmath}
\usepackage{amssymb}
\usepackage{stmaryrd}

% this should be the last package used
\usepackage{pdfsetup}

\urlstyle{rm}
\isabellestyle{it}

\begin{document}

\title{Wiener Measure as a Projective Limit}
\author{Dmitriy Traytel}
\maketitle

\begin{abstract}
  Brownian motion is the basic object of continuous-time stochastic analysis.
  This entry constructs it in Isabelle, with the Archive entry
  \texttt{Kolmogorov\_Chentsov} supplying the criterion that turns the
  construction into one with continuous paths.  We build Wiener measure as
  the projective limit of finite-dimensional distributions assembled from
  Gaussian increments, prove those increments independent and Gaussian, and
  compute the fourth moment
  $\mathbb{E}\,|B_t - B_s|^{4} = 3\,(t-s)^{2}$, which is exactly the
  hypothesis the Kolmogorov--Chentsov theorem asks for.  Projectivity reduces
  to the convolution law for Gaussian densities by merging two adjacent
  increments, which brings the construction within reach of the
  Daniell--Kolmogorov theorem already available in \texttt{HOL-Probability}.
  We follow the account of Karatzas and Shreve~\cite{KaratzasShreve}.
\end{abstract}

\tableofcontents

\parindent 0pt\parskip 0.5ex

\input{session}

\bibliographystyle{abbrv}
\bibliography{root}

\end{document}
